\documentclass{article}

% Ćirilica: kodiranje fonta i jezik (srpski, ćirilica)
\usepackage[T2A]{fontenc}
\usepackage[utf8]{inputenc}
\usepackage[serbianc]{babel}

% Veličina strane i margine
\usepackage[a4paper,top=2cm,bottom=2cm,left=3cm,right=3cm,marginparwidth=1.75cm]{geometry}

\usepackage[european]{circuitikz}

% Korisni paketi
\usepackage{amsmath}
\usepackage{graphicx}
\usepackage{float}
\usepackage{booktabs}
\usepackage{tcolorbox}
\tcbuselibrary{breakable}
\usepackage[colorlinks=true, allcolors=blue]{hyperref}

% Цртежи (TikZ)
\usepackage{tikz}
\usetikzlibrary{arrows.meta, calc, decorations.markings, babel}

% TikZ стилови који се понављају на цртежима
\tikzset{
    naboj/.style={circle, draw, thick, minimum size=6mm, inner sep=0pt, font=\small},
    jezgro/.style={circle, draw, thick, fill=gray!15, minimum size=7mm, inner sep=0pt, font=\small},
    elektron/.style={circle, draw, fill=cyan!30, inner sep=0pt, minimum size=5pt},
    supljina/.style={circle, draw, thick, minimum size=5mm, inner sep=0pt, font=\small},
    sredstrelica/.style={postaction={decorate}, decoration={markings, mark=at position 0.55 with {\arrow{Stealth}}}},
    % --- графици наизменичних величина ---
    pomocna/.style={dashed, gray!70, thin},
    mreza/.style={gray!25, very thin},
    kriva/.style={very thick, red!75!black},
    % --- проводник у пресеку: празан, тачка (ка нама), крстић (од нас) ---
    provodnik/.style={circle, draw, thick, fill=white, minimum size=8pt, inner sep=0pt},
    provodnikT/.style={provodnik, path picture={\fill[black] (path picture bounding box.center) circle (1.1pt);}},
    provodnikK/.style={provodnik, path picture={
        \draw (path picture bounding box.south west) -- (path picture bounding box.north east)
              (path picture bounding box.north west) -- (path picture bounding box.south east);}},
    oznaka/.style={fill=white, inner sep=1pt, font=\small},
}

% Оквир за решене примере
\newtcolorbox[auto counter]{primer}{breakable,
    colback=white, colframe=black!55, colbacktitle=black!55, coltitle=white,
    fonttitle=\bfseries\small, title={Пример~\thetcbcounter},
    boxrule=0.5pt, arc=1.5pt, left=6pt, right=6pt, top=3pt, bottom=3pt,
    before skip=8pt, after skip=10pt}

% Оквир за кључне формуле
\newtcolorbox{kljucno}{colback=blue!4, colframe=blue!45!black,
    boxrule=0.5pt, arc=1.5pt, left=8pt, right=8pt, top=4pt, bottom=4pt}

\title{ОЕТ 1 - Вежбање}
\author{Саобраћајно техничка школа, Београд, Земун}

\begin{document}
\maketitle


\begin{primer}
У колу на слици извор има електромоторну силу $E = 12\,\text{V}$, а
отпорници имају отпорности $R_1 = 2\,\Omega$ и $R_2 = 4\,\Omega$. Израчунати:
\begin{enumerate}
    \item[а)] еквивалентну отпорност кола,
    \item[б)] јачину струје у колу,
    \item[в)] напоне на отпорницима $R_1$ и $R_2$.
\end{enumerate}

\begin{center}
\begin{circuitikz}
    \draw (0,0) to[battery1, l=$E$, invert] (0,3)
          to[short, i>^=$I$] (1,3)
          to[R, l_=$R_1$, v^=$U_1$] (4,3)
          to[R, l_=$R_2$, v^=$U_2$] (4,0)
          -- (0,0);
\end{circuitikz}
\end{center}

\textbf{Решење:}

а) Отпорници су везани редно, па се њихове отпорности сабирају:
\[
R_e = R_1 + R_2 = 2\,\Omega + 4\,\Omega = 6\,\Omega
\]

б) Струју рачунамо из Омовог закона:
\[
I = \frac{E}{R_e} = \frac{12\,\text{V}}{6\,\Omega} = 2\,\text{A}
\]

в) Кроз оба отпорника тече иста струја $I$, па су напони:
\[
U_1 = I \cdot R_1 = 2\,\text{A} \cdot 2\,\Omega = 4\,\text{V}, \qquad
U_2 = I \cdot R_2 = 2\,\text{A} \cdot 4\,\Omega = 8\,\text{V}
\]

\textbf{Провера:} збир напона на отпорницима мора бити једнак електромоторној
сили извора (други Кирхофов закон):
\[
U_1 + U_2 = 4\,\text{V} + 8\,\text{V} = 12\,\text{V} = E
\]
\end{primer}

\begin{primer}
Два отпорника, $R_1 = 6\,\Omega$ и $R_2 = 3\,\Omega$, везана су паралелно и
прикључена на извор електромоторне силе $E = 12\,\text{V}$. Унутрашњу
отпорност извора занемарити. Израчунати:
\begin{enumerate}
    \item[а)] еквивалентну отпорност кола,
    \item[б)] струју коју даје извор,
    \item[в)] струје кроз отпорнике $R_1$ и $R_2$.
\end{enumerate}

\begin{center}
\begin{circuitikz}
    \draw (0,0) to[battery1, l=$E$, invert] (0,3)
          to[short, i>^=$I$, -*] (3,3)
          to[R, l_=$R_1$, i>^=$I_1$, -*] (3,0)
          -- (0,0);
    \draw (3,3) -- (6,3)
          to[R, l_=$R_2$, i>^=$I_2$] (6,0)
          -- (3,0);
\end{circuitikz}
\end{center}

\textbf{Решење:}

а) Отпорници су везани паралелно, па је за два отпорника:
\[
R_e = \frac{R_1 R_2}{R_1 + R_2} = \frac{6\,\Omega \cdot 3\,\Omega}{6\,\Omega + 3\,\Omega}
    = \frac{18}{9}\,\Omega = 2\,\Omega
\]

б) Струја извора, из Омовог закона:
\[
I = \frac{E}{R_e} = \frac{12\,\text{V}}{2\,\Omega} = 6\,\text{A}
\]

в) Оба отпорника су прикључена директно на извор, па је на сваком напон
$U = E = 12\,\text{V}$:
\[
I_1 = \frac{U}{R_1} = \frac{12\,\text{V}}{6\,\Omega} = 2\,\text{A}, \qquad
I_2 = \frac{U}{R_2} = \frac{12\,\text{V}}{3\,\Omega} = 4\,\text{A}
\]

\textbf{Провера:} збир струја кроз отпорнике мора бити једнак струји извора
(први Кирхофов закон):
\[
I_1 + I_2 = 2\,\text{A} + 4\,\text{A} = 6\,\text{A} = I
\]
\end{primer}

\begin{primer}
Бакарни проводник има дужину $l = 100\,\text{m}$ и површину попречног
пресека $S = 2{,}5\,\text{mm}^2$. Специфична отпорност бакра је
$\rho = 0{,}0175\,\Omega\,\dfrac{\text{mm}^2}{\text{m}}$.
Израчунати отпорност проводника.

\textbf{Решење:}

Отпорност проводника рачунамо по формули:
\[
R = \rho \, \frac{l}{S}
  = 0{,}0175\,\Omega\,\frac{\text{mm}^2}{\text{m}} \cdot \frac{100\,\text{m}}{2{,}5\,\text{mm}^2}
  = 0{,}0175 \cdot 40\,\Omega
  = 0{,}7\,\Omega
\]

Јединице $\text{mm}^2$ и $\text{m}$ се скраћују, па остаје само $\Omega$.
\end{primer}

\begin{primer}
Кондензатор је прикључен на напон $U = 12\,\text{V}$ и на њему се нагомилало
наелектрисање $Q = 60\,\mu\text{C}$. Израчунати:
\begin{enumerate}
    \item[а)] капацитивност кондензатора,
    \item[б)] наелектрисање кондензатора ако се напон повећа на $U' = 24\,\text{V}$.
\end{enumerate}

\textbf{Решење:}

а) Капацитивност рачунамо по формули:
\[
C = \frac{Q}{U} = \frac{60\,\mu\text{C}}{12\,\text{V}} = 5\,\mu\text{F}
\]

б) Капацитивност зависи само од грађе кондензатора, па се не мења када се
промени напон. Ново наелектрисање је:
\[
Q' = C \cdot U' = 5\,\mu\text{F} \cdot 24\,\text{V} = 120\,\mu\text{C}
\]

Када се напон удвостручи, удвостручи се и наелектрисање, а капацитивност
остаје иста.
\end{primer}

%==================================================================
% БЕЛЕШКЕ ИЗ СВЕСКЕ (редом по сликама IMG_6501 -- IMG_6508)
%==================================================================

\end{document}